\section{Related Work}
\label{sec:related}
%% ============================================================

The economic case for ISRU has been examined primarily for mission-specific applications. Sanders and Larson \cite{sanders2015} found lunar oxygen extraction cost-effective after approximately five crewed missions, but their model does not generalize to manufactured goods. Crawford \cite{crawford2015} catalogued lunar mineralogical inventories; Cilliers et al.\ \cite{cilliers2023} examined regolith processing pathways; Kornuta et al.\ \cite{kornuta2019} developed a commercial lunar propellant architecture. The MOXIE experiment \cite{hecht2021} provided empirical validation of in-situ oxygen production on Mars, though at sub-manufacturing scales. Sowers \cite{sowers2021,sowers2023} presented NPV-based business cases for lunar ice mining and cislunar economics. Metzger et al.\ \cite{metzger2013} proposed a bootstrapping approach in which a seed factory progressively builds its own capacity. Asteroid mining NPV frameworks have been developed by Sonter \cite{sonter1997}, Elvis \cite{elvis2012}, and Andrews et al.\ \cite{andrews2015}; Ishimatsu et al.\ \cite{ishimatsu2016} modeled Earth--Moon--Mars logistics as a multicommodity network flow. These analyses focus on resource extraction rather than structural manufacturing. Lunar additive manufacturing demonstrations \cite{werkheiser2015,cesaretti2014} provide empirical grounding for the ISRU manufacturing pathway assumed here. O'Neill \cite{oneill1974,oneill1976} argued on energetic grounds for lunar material delivery to orbit, but lacked a parametric cost model and did not incorporate learning curves. The real options framework \cite{dixit1994} extends NPV for irreversible investments under uncertainty; Saleh et al.\ \cite{saleh2003} applied flexibility analysis to spacecraft, and de~Weck et al.\ \cite{deweck2004} demonstrated staged deployment value. Organizational forgetting \cite{benkard2000,thompson2012} has been documented in aerospace production lines.

Jones \cite{jones2018,jones2020,jones2022} documented a roughly two-orders-of-magnitude reduction in launch cost per kilogram over six decades. A critical distinction is between vehicle manufacturing cost (which follows a learning curve) and per-kilogram delivery cost (dominated by propellant and operations, largely independent of cumulative launches). Falcon~9 reuse economics \cite{zapata2019} confirm that per-launch cost reductions arise from fixed-cost amortization at high flight rates, not production learning. Wertz \cite{wertz2011} estimated launch vehicle production learning rates of 85--90\%; the NASA Cost Estimating Handbook \cite{nasa2015handbook} provides standardized parametric methodologies that inform our parameter selection. A recent analysis of NASA's internal cost database \cite{kim2025} finds launch cost \emph{per kilogram paid by NASA} rising by an average of 2.8\% annually from 1996--2024 even after new entrants joined the market---a reminder that our declining-cost assumption reflects commercial list prices and vehicle-level learning, not government procurement cost, and that the crossover timing is sensitive to which cost trend is assumed. Our model adopts a two-component launch cost: a propellant operational asymptote plus a learnable component ($\mathrm{LR}_L = 0.97$); we sweep $\mathrm{LR}_L$ from 0.90 to 1.00 in \S\ref{sec:sensitivity}.

The Wright learning curve \cite{wright1936}---in which unit cost decreases as a power function of cumulative production---is well established in aerospace economics \cite{dutton1984,argote1990,nagy2013}; we formalize it in \S\ref{sec:model} (Eq.~\ref{eq:earth_mfg}). Typical aerospace learning rates range from 0.80 to 0.95 \cite{wertz2011}; Baumers et al.\ \cite{baumers2016} documented 0.85--0.92 for metal additive manufacturing, and Ding et al.\ \cite{ding2021} generalize the additive-manufacturing cost model to explicitly account for process-failure rates, a cost driver our ISRU manufacturing-yield parameter is intended to capture in reduced form. We are not aware of prior work applying the Wright model to a comparative Earth-versus-ISRU manufacturing analysis with NPV timing.

%% ============================================================
